% Section 18 Header
\section{Pattern 16: Greedy Algorithms}

\begin{frame}[c]{\empty}
    \begin{center}
        \Large\textbf{\secname}
    \end{center}
\end{frame}

% Slide 1: Concept
\begin{frame}{What are Greedy Algorithms?}
  \begin{itemize}
    \item \textbf{The Core Idea:} An algorithmic paradigm that builds up a solution piece by piece, always choosing the next piece that offers the most obvious and immediate local benefit ("greedy" choice).
    \item \textbf{Why it works:} For certain problems, making locally optimal decisions leads to a globally optimal solution.
    \item \textbf{Tradeoffs:}
      \begin{itemize}
        \item Highly efficient (usually $O(N)$ or $O(N \log N)$).
        \item Difficult to mathematically prove correctness.
        \item If greedy fails, you must fall back to Dynamic Programming (which checks all combinations).
      \end{itemize}
  \end{itemize}
\end{frame}

% Problem 1: Gas Station
\begin{frame}[fragile]{Problem: Circular Route Fuel Stations}
  \textbf{Scenario:}
  Given `fuel` and `cost` arrays for circular route fuel stations, find the starting station index from which you can travel around the circuit once without running out of gas.
  
  \vspace{0.3cm}
  \textbf{Examples:}
  \begin{itemize}
    \item \textbf{Input:} \texttt{fuel = [1,2,3,4,5], cost = [3,4,5,1,2]} $\rightarrow$ \textbf{Output:} \texttt{3}
  \end{itemize}
\end{frame}

% Problem 2: Candy
\begin{frame}[fragile]{Problem: Min Candy Distribution}
  \textbf{Scenario:}
  There are n children standing in a line. Each child has a rating. Distribute candies such that each child has at least 1, and children with higher ratings than neighbors get more.
  
  \vspace{0.3cm}
  \textbf{Examples:}
  \begin{itemize}
    \item \textbf{Input:} \texttt{[1,0,2]} $\rightarrow$ \textbf{Output:} \texttt{5}
  \end{itemize}
\end{frame}

% Problem 3: Queue Reconstruction by Height
\begin{frame}[fragile]{Problem: Line Reconstruction by Height Attributes}
  \textbf{Scenario:}
  Reconstruct a queue where each person is represented by [h, k] where h is their height and k is the number of people in front of them who are taller or equal.
  
  \vspace{0.3cm}
  \textbf{Examples:}
  \begin{itemize}
    \item \textbf{Input:} \texttt{[[7,0],[4,4],[7,1],[5,0],[6,1],[5,2]]} $\rightarrow$ \textbf{Output:} \texttt{[[5,0],[7,0],[5,2],[6,1],[4,4],[7,1]]}
  \end{itemize}
\end{frame}

% Common Triggers Slide
\begin{frame}{\secname\ -- Common Triggers}
  \begin{itemize}
    \item Scheduling, interval selection, minimizations.
        \item Gas stations, container loading, coin distributions (canonical cases).
  \end{itemize}
\end{frame}
